\section{Results and System Interpretation}
\subsection{STRICT Cross-Sample Parallelism}
On the prior eight-core GEMM cohort, STRICT Pipeline reduces construction from 149.56 to 20.95 seconds and achieves $1.638\pm0.005\times$ paired Process speedup over Stock; Group takes 37.88 seconds. Construction is about 45\% of Stock Process. Matched Group/Flat/Pipeline controls isolate C1 from other Stock repairs.

\begin{figure}[t]\centering\begin{tikzpicture}
\begin{groupplot}[group style={group size=2 by 1,horizontal sep=0.8cm},width=4.15cm,height=3.7cm,xtick={8,16},xmin=7,xmax=17,ymin=0,ymax=40,xlabel={Physical cores},ylabel={Build (s)},label style={font=\scriptsize},tick label style={font=\scriptsize},legend style={font=\scriptsize,at={(1.15,1.08)},anchor=south,legend columns=3},error bars/y dir=both,error bars/y explicit]
\nextgroupplot[title={},ylabel={Build (s)}]
\addplot+[mark=square*] coordinates {(8,17.413333333333334) +- (0,0.3250128202599597) (16,17.683333333333334) +- (0,0.3257811126098835)};
\addplot+[mark=triangle*] coordinates {(8,9.32) +- (0,0.15394804318340596) (16,5.506666666666667) +- (0,0.030550504633038766)};
\addplot+[mark=o] coordinates {(8,9.363333333333333) +- (0,0.10598742063723142) (16,5.3933333333333335) +- (0,0.08082903768654771)};
\legend{Group,Flat,Pipeline}
\nextgroupplot[title={},ylabel={}]
\addplot+[mark=square*] coordinates {(8,34.82333333333333) +- (0,0.3181718613160717) (16,33.99) +- (0,0.6602272336097627)};
\addplot+[mark=triangle*] coordinates {(8,19.553333333333335) +- (0,0.6217180497084931) (16,10.896666666666667) +- (0,0.040414518843273746)};
\addplot+[mark=o] coordinates {(8,18.993333333333332) +- (0,0.11503622617824955) (16,10.846666666666668) +- (0,0.04041451884327389)};
\end{groupplot}
\end{tikzpicture}

\caption{Construction means and sample SD, three processes per arm/point. Same-cohort Group/Flat/Pipeline controls on BNN (left) and DLA (right); capped routing tails do not establish full-route legality. Error bars are small at this scale.}\label{fig:causal}\end{figure}

With ten groups of four samples, Group limits concurrency to four searches. At eight cores, BNN Group/Flat/Pipeline take $17.41\pm0.33$, $9.32\pm0.15$ and $9.36\pm0.11$ s (Fig.~\ref{fig:causal}). At sixteen cores the times are 17.68/5.51/5.39 s. DLA shows the same pattern (34.82/19.55/18.99 s at eight cores). Both global schedulers benefit from exposing work across groups. Separate eight-core BNN-2D controls take 7.75/4.23/4.16 s on a different graph.

\subsection{Wave Scheduling: Concurrency Versus Work}
\begin{table}[t]\centering\caption{Historical BNN40 bucket comparison. Work profiles and warm complete-batch seconds on eight physical cores and one H100.}\label{tab:waves}
\begin{tabular}{lrr}\toprule Metric & Shift 0 & Shift 20\\\midrule
\tableinput{tables/waves.tex}\midrule\tableinput{tables/batch_seconds.tex}\bottomrule\end{tabular}\end{table}
Table~\ref{tab:waves} links independently profiled work to warm complete-batch service on the same BNN40 input. Coarsening the ALT bucket reduces total waves $9.112\times$ and maximum sample waves $10.627\times$. Eligible edge work rises $0.3504\%$, while edge-weighted frontier grows $6.6484\times$. The eight-core CPU gains \BatchCpuGain\ and the GPU gains \BatchCudaGain. Wider fronts expose more simultaneous work between synchronization points. Both shifts retain FP32 labels and agree on all sample outputs on this batch.

Service includes reset, search, closure and materializing outputs, including GPU readback and synchronization. Load, preparation and verification fall outside the timer. Each time is the median of three process medians, each over three warm batches with profiling and witnesses off. Speedups here divide those medians. Native E2 additionally charges graph conversion and outer lookahead construction.

\subsection{Matched Native Process and Construction}
\begin{table*}[t]\centering\caption{Complete native Process seconds and directly paired speedups, mean $\pm$ sample SD ($n=3$). S: Stock; P: Pipeline; C: ALT CUDA. GEMM/BNN are one new allocation; LSTM is a separate historical cohort.}\label{tab:matchedproc}
\begin{tabular}{lrrrrrr}\toprule Workload & Stock (s) & Pipeline (s) & ALT CUDA (s) & S/P & P/C & S/C\\\midrule
\tableinput{tables/matched_process.tex}\bottomrule\end{tabular}\end{table*}
\begin{table*}[t]\centering\caption{Native lookahead construction E2 seconds and directly paired speedups for the same runs as Table~\ref{tab:matchedproc}; mean $\pm$ sample SD, $n=3$.}\label{tab:matchede2}
\begin{tabular}{lrrrrrr}\toprule Workload & Stock (s) & Pipeline (s) & ALT CUDA (s) & S/P & P/C & S/C\\\midrule
\tableinput{tables/matched_e2.tex}\bottomrule\end{tabular}\end{table*}
ALT CUDA achieves E2 speedups of \GemmNewPipeBuild\ for GEMM and \BnnNewPipeBuild\ for BNN over Pipeline (Tables~\ref{tab:matchedproc}--\ref{tab:matchede2}). E2 accounts for \EstimateGemmShare\ and \EstimateBnnShare\ of mean Pipeline Process time. Holding the mean remainder fixed, eliminating E2 would give \EstimateGemmZero\ and \EstimateBnnZero\ Process speedups. Substituting mean CUDA E2 instead gives \EstimateGemmReplaced\ for each workload. Actual paired Process speedups are \GemmNewPipeProc\ and \BnnNewPipeProc; other stages can also change.


\subsection{Matched Native Engine and Backend Comparison}
\begin{table}[t]\centering
\caption{Independent eight-core/H100 cohort. Directly paired speedups ($n=3$, mean $\pm$ sample SD). P: Pipeline, A: ALT CPU, C: ALT CUDA.}\label{tab:engines}
\begin{tabular}{@{}lrrr@{}}\toprule Case and scope & P/A & A/C & P/C\\\midrule
\tableinput{tables/engine_ratios.tex}\bottomrule\end{tabular}\end{table}
Changing the CPU engine and scheduler together yields Process ratios near one (Table~\ref{tab:engines}), including $0.973\pm0.074$ on GEMM. Holding the ALT rule fixed, CUDA speeds construction by $4.690\times$ on GEMM and $4.331\times$ on BNN. The corresponding Process speedups are $1.146\times$ and $1.162\times$.

\subsection{Native LSTM: Five Matched Configurations}
\begin{table}[t]\centering\caption{Historical LSTM native seconds, mean $\pm$ sample SD ($n=3$); five arms on its own rotated eight-core/H100 cohort.}\label{tab:lstm}
\begin{tabular}{lrr}\toprule Method & E2 & Process\\\midrule
\tableinput{tables/lstm.tex}\bottomrule\end{tabular}\end{table}
On historical LSTM, CUDA achieves \LstmPipeBuild\ E2 and \LstmPipeProc\ Process speedups over Pipeline. Against matched ALT CPU the speedups are \LstmAltCpuBuild\ and \LstmAltCpuProc. Table~\ref{tab:lstm} includes the full five-arm comparison.

\subsection{Semantics, QoR and Size Dependence}
\begin{table}[t]\centering\caption{Native qualification: nodes/edges in millions; CPD/WL changes (\%) against Stock; $\Delta q$ counts final-map congestion-word changes against STRICT. All cases have two layers and 40 samples; LSTM is historical.}\label{tab:qor}
\begin{tabular}{lrrrrr}\toprule Case & $N/10^6$ & $M/10^6$ & CPD & WL & $\Delta q$\\\midrule
\tableinput{tables/qualification_compact.tex}\bottomrule\end{tabular}\end{table}
Complete ALT CPU/CUDA samples and maps agree. Stock/Pipeline maps also agree. ALT changes only congestion representatives (Table~\ref{tab:qor}), with zero delay or missing-state differences. Each route passes independent FULL/STA with zero overuse and preserves its own SHA. All native runs meet their respective CPD/WL gates.

On the historical BNN cohort, CUDA achieves $3.274\pm0.222\times$ E2 and $1.091\pm0.011\times$ Process speedups over Flat. On small s38417 the corresponding ratios fall to $0.274\pm0.032\times$ and $0.912\pm0.016\times$. Initialization, scans and launches can outweigh search savings.

\subsection{Closest Alternative and Cache Applicability}\label{sec:applicability}
The locked \texttt{compressed\_map} requires one layer, excluding these two-layer graphs. Separable remains unmeasured. Ordinary caching supports valid reuse: 36 fixed-input cache hits passed independent FULL/STA. Our timings cover fresh construction.
